\documentclass{resume} % Use the custom resume.cls style
\usepackage[T1,T2A]{fontenc}
\usepackage[utf8]{inputenc}

\usepackage{cmap} 
\usepackage{fontawesome}  % для значков в шапке
\usepackage[english,bulgarian,ukrainian,russian]{babel}
\usepackage[left=0.4in,top=0.3in,right=0.4in,bottom=0.3in]{geometry} % Document margins
\newcommand{\tab}[1]{\hspace{.2667\textwidth}\rlap{#1}} 
\newcommand{\itab}[1]{\hspace{0em}\rlap{#1}}
\name{Кирякин Максим} % Your name
% You can merge both of these into a single line, if you do not have a website.
\address{+7(916) 344-6061 \\ Россия, Москва } 
\address{
	\faEnvelope \space \href{mailto:makxim.kiryakin@gmail.com}{makxim.kiryakin@gmail.com}  \faGithub \space \href{https://github.com/MaximKiryakin}{www.github.com} 
	\faPaperPlane \space \href{https://t.me/Maxim_Kiryakin}{@Maxim\_Kiryakin}
}  %

\def\sectionskip{\smallskip} % tighter spacing between sections to fit one page
\def\sectionlineskip{\smallskip} % tighter space above each section rule
\linespread{0.97} % shave line spacing slightly to keep everything on one page

\begin{document}

	\begin{rSection}{Кратко}
		Data Scientist, выводящий ML в продакшн на больших данных. Построил пайплайн кластеризации
		\textbf{120+ млн клиентских профилей} на PySpark и вывел NPV-модели в промышленный контур одного
		из крупнейших банков страны --- они влияют на реальные решения по продуктам. Магистр ВМК МГУ.
	\end{rSection}

	\begin{rSection}{Образование}
		\textbf{МГУ имени М.В. Ломоносова (Магистратура)} \hfill {Сентябрь 2024 - Настоящее время}\\
		Факультет вычислительной математики и кибернетики (\textit{Средний балл:} 4.5).\\
		\textbf{МГУ имени М.В. Ломоносова (Бакалавриат)} \hfill {Сентябрь 2020 - июнь 2024}\\
		Факультет вычислительной математики и кибернетики (\textit{Средний балл:} 4.3).
	\end{rSection}

	\begin{rSection}{Опыт работы}
		\textbf{Сбер} (топ-3 банка страны по активам, 100+ млн клиентов) \hfill {Май 2024 - Настоящее время}\\
		Data Scientist в Дирекции Инвестиционной Экспертизы (Блок Финансы) \\
		\vspace{-1.25em}
		\begin{itemize}
			\itemsep -3pt {}
			\item Построил и внедрил пайплайн кластеризации \textbf{120+ млн клиентских профилей} на PySpark (Hadoop); автоматически пересчитывается [\textbf{еженедельно}], заменив ручную сегментацию.
			\item Вывел ML- и эконометрические NPV-модели в промышленный контур Банка, повысив точность оценки на \textbf{[X]\%}; они влияют на решения по ценообразованию. Выстроил мониторинг, переобучение и бизнес-приёмку.
			\item Спроектировал и пилотировал новую LTV-метрику эффективности продуктов: определил с нуля и подтвердил статистическими экспериментами до внедрения в дирекции.
			\item Стек: Python (PySpark, Pandas, Scikit-learn), PyTorch, SQL, Hive, Git.
		\end{itemize}
			
		\textbf{Сбер} \hfill {Август 2023 - Март 2024}\\
		Стажёр-Data Scientist в Управлении структуры баланса (Блок Финансы) \\
		\vspace{-1.25em}
		\begin{itemize}
			\itemsep -3pt {}
			\item Внедрил предобученную ML-модель в продуктовый пайплайн, снизив ошибку прогноза целевой переменной (MAE) на \textbf{30\%} на валидации при воспроизводимом инференсе.
			\item Ускорил ключевой расчёт в \textbf{3 раза} и снизил потребление памяти за счёт профилирования, алгоритмического рефакторинга и векторизации.
			\item Построил автоматические PySpark ETL-процессы выгрузки и обработки больших данных из HDFS с логированием и обработкой ошибок, заменив хрупкие ручные прототипы; покрыл тестами и документацией.
			\item Стек: Python (PySpark, Pandas, Scikit-learn), SQL, Hive.
		\end{itemize}
	\end{rSection}

	\begin{rSection}{Навыки}
		\begin{tabular}{ @{} >{\bfseries}l @{\hspace{6ex}} l }
			Языки: & Python, SQL, C/C++ \\
			ML и данные: & PySpark, Scikit-learn, PyTorch, Pandas, NumPy, SciPy, Hive, Hadoop \\
			Методы ML: & Кластеризация, временные ряды, feature engineering, эконометрика, A/B-тесты \\
			Продакшн: & End-to-end пайплайны, деплой моделей, мониторинг и переобучение, Git
		\end{tabular}\\
	\end{rSection}

	\begin{rSection}{Повышение квалификации}
		\begin{itemize}
			\itemsep -3pt {}
			\item <<Методы анализа данных и машинного обучения>> --- Сбер Университет
			\hfill \href{https://disk.yandex.ru/i/7NCFgcRKg5lyYg}{Сертификат (с отличием)}
			\item <<Введение в финансовую математику>>, <<Финансовая эконометрика>> --- Фонд <<Институт ``Вега''>>
			\item <<Математические модели в инвестиционных банках>> --- Технологический Центр Дойче Банка
		\end{itemize}
	\end{rSection}

	\begin{rSection}{Достижения}
		\begin{itemize}
			\itemsep -3pt {}
			\item \textbf{Олимпиада ``Шаг в будущее''} --- диплом 3 степени, топ-$60$ среди $1600+$ участников.
			\hfill \href{https://disk.yandex.ru/i/43o4Hafel2AneA}{Диплом}
		\end{itemize}
	\end{rSection}

	\begin{rSection}{О себе}
		\begin{itemize}
			\itemsep -3pt {}
			\item \textbf{Языки}: Русский (родной), Английский (продвинутый). \quad \textbf{Спорт}: 9 лет самбо, призёр первенства Берлина.
		\end{itemize}
	\end{rSection}

		

\end{document}
